\documentclass[UTF8,12pt]{ctexart}
\usepackage[a4paper,margin=24mm]{geometry}
\usepackage{amsmath,amssymb,bm,graphicx,adjustbox}
\newcommand{\fitmath}[1]{\begin{adjustbox}{max width=\linewidth}$\displaystyle #1$\end{adjustbox}}
\setlength{\parindent}{0pt}
\setlength{\parskip}{.7em}
\sloppy
\allowdisplaybreaks
\begin{document}
\section*{2019 年考研数学三 · 真题}
\subsection*{原 PDF 第 1 页}

\subsection*{一、选择题（本题共 8 小题，每小题 4 分，共 32 分。在每小题给出的四个选项中，只有一项符合题目要求，把所选项前的字母填在题后的括号内。）}

（1）当 \(\displaystyle x\to0\) 时，若 \(\displaystyle x-\tan x\) 与 \(\displaystyle x^k\) 是同阶无穷小，则 \(\displaystyle k=(\quad)\)。


（A）1。　（B）2。　（C）3。　（D）4。


（2）已知方程 \(\displaystyle x^5-5x+k=0\) 有 3 个不同的实根，则 \(\displaystyle k\) 的取值范围是（　）。


（A）\(\displaystyle (-\infty,\allowbreak -4)\)。　（B）\(\displaystyle (4,\allowbreak +\infty)\)。　（C）\(\displaystyle \{-4,\allowbreak 4\}\)。　（D）\(\displaystyle (-4,\allowbreak 4)\)。


（3）已知微分方程 \(\displaystyle y^{\prime\prime}+ay^\prime+by=ce^x\) 的通解为 \(\displaystyle y=(C_1+C_2x)e^{-x}+e^x\)，则 \(\displaystyle a,\allowbreak b,\allowbreak c\) 依次为（　）。


（A）1，0，1。　（B）1，0，2。　（C）2，1，3。　（D）2，1，4。


（4）若 \(\displaystyle \sum_{n=1}^{\infty}nu_n\) 绝对收敛，\(\displaystyle \sum_{n=1}^{\infty}\dfrac{\displaystyle v_n}{\displaystyle n}\) 条件收敛，则（　）。


（A）\(\displaystyle \sum_{n=1}^{\infty}u_nv_n\) 条件收敛。　（B）\(\displaystyle \sum_{n=1}^{\infty}u_nv_n\) 绝对收敛。


（C）\(\displaystyle \sum_{n=1}^{\infty}(u_n+v_n)\) 收敛。　（D）\(\displaystyle \sum_{n=1}^{\infty}(u_n+v_n)\) 发散。


（5）设 \(\displaystyle \boldsymbol A\) 是 4 阶矩阵，\(\displaystyle \boldsymbol A^*\) 是 \(\displaystyle \boldsymbol A\) 的伴随矩阵，若线性方程组 \(\displaystyle \boldsymbol A\boldsymbol x=\boldsymbol0\) 的基础解系中只有 2 个向量，则 \(\displaystyle r(\boldsymbol A^*)=(\quad)\)。


（A）0。　（B）1。　（C）2。　（D）3。


（6）设 \(\displaystyle \boldsymbol A\) 是 3 阶实对称矩阵，\(\displaystyle \boldsymbol E\) 是 3 阶单位矩阵。若 \(\displaystyle \boldsymbol A^2+\boldsymbol A=2\boldsymbol E\)，且 \(\displaystyle |\boldsymbol A|=4\)，则二次型 \(\displaystyle \boldsymbol x^{\mathrm T}\boldsymbol A\boldsymbol x\) 的规范形为（　）。


（A）\(\displaystyle y_1^2+y_2^2+y_3^2\)。　（B）\(\displaystyle y_1^2+y_2^2-y_3^2\)。　（C）\(\displaystyle y_1^2-y_2^2-y_3^2\)。　（D）\(\displaystyle -y_1^2-y_2^2-y_3^2\)。


（7）设 \(\displaystyle A,\allowbreak B\) 为随机事件，则 \(\displaystyle P(A)=P(B)\) 的充分必要条件是（　）。


（A）\(\displaystyle P(A\cup B)=P(A)+P(B)\)。　（B）\(\displaystyle P(AB)=P(A)P(B)\)。


（C）\(\displaystyle P(A\overline B)=P(B\overline A)\)。　（D）\(\displaystyle P(AB)=P(\overline A\,\overline B)\)。


（8）设随机变量 \(\displaystyle X\) 与 \(\displaystyle Y\) 相互独立，且都服从正态分布 \(\displaystyle N(\mu,\allowbreak \sigma^2)\)，则 \(\displaystyle P\{|X-Y|<1\}\)（　）。


（A）与 \(\displaystyle \mu\) 无关，而与 \(\displaystyle \sigma^2\) 有关。　（B）与 \(\displaystyle \mu\) 有关，而与 \(\displaystyle \sigma^2\) 无关。


（C）与 \(\displaystyle \mu,\allowbreak \sigma^2\) 都有关。　（D）与 \(\displaystyle \mu,\allowbreak \sigma^2\) 都无关。


\subsection*{二、填空题（本题共 6 小题，每小题 4 分，共 24 分。把答案填在题中横线上。）}

（9）\(\displaystyle \lim_{n\to\infty}\left[\dfrac{\displaystyle 1}{\displaystyle 1\cdot2}+\dfrac{\displaystyle 1}{\displaystyle 2\cdot3}+\cdots+\dfrac{\displaystyle 1}{\displaystyle n(n+1)}\right]^n=\underline{\qquad}\)。


（10）曲线 \(\displaystyle y=x\sin x+2\cos x\;(-\dfrac{\displaystyle \pi}{\displaystyle 2}<x<\dfrac{\displaystyle 3\pi}{\displaystyle 2})\) 的拐点坐标为 \(\displaystyle \underline{\qquad}\)。


（11）已知函数 \(\displaystyle f(x)=\displaystyle \int_1^x\sqrt{1+t^4}\,\mathrm dt\)，则 \(\displaystyle \int_0^1x^2f(x)\,\mathrm dx=\underline{\qquad}\)。


\clearpage

\subsection*{原 PDF 第 2 页}

（12）以 \(\displaystyle P_A,\allowbreak P_B\) 分别表示 A、B 两个商品的价格，设商品 A 的需求函数 \(\displaystyle Q_A=500-P_A^2-P_AP_B+2P_B^2\)，则当 \(\displaystyle P_A=10,\allowbreak P_B=20\) 时，商品 A 的需求量对自身价格的弹性 \(\displaystyle \eta_{AA}\;(\eta_{AA}>0)=\underline{\qquad}\)。


（13）已知矩阵 \(\displaystyle \boldsymbol A=\begin{pmatrix}1&0&-1\\1&1&-1\\0&1&a^2-1\end{pmatrix}\)，\(\displaystyle \boldsymbol b=\begin{pmatrix}0\\1\\a\end{pmatrix}\)，若线性方程组 \(\displaystyle \boldsymbol A\boldsymbol x=\boldsymbol b\) 有无穷多解，则 \(\displaystyle a=\underline{\qquad}\)。


（14）设随机变量 \(\displaystyle X\) 的概率密度为 \(\displaystyle f(x)=\begin{cases}\dfrac{\displaystyle x}{\displaystyle 2},\allowbreak &0<x<2,\allowbreak \\0,\allowbreak &\text{其他},\allowbreak \end{cases}\) \(\displaystyle F(x)\) 为 \(\displaystyle X\) 的分布函数，\(\displaystyle E(X)\) 为 \(\displaystyle X\) 的数学期望，则 \(\displaystyle P\{F(X)>E(X)-1\}=\underline{\qquad}\)。


\subsection*{三、解答题（本题共 9 小题，共 94 分。解答应写出文字说明、证明过程或演算步骤。）}

（15）（本题满分 10 分）已知函数 \(\displaystyle f(x)=\begin{cases}x^{2x},\allowbreak &x>0,\allowbreak \\xe^x+1,\allowbreak &x\leq0.\end{cases}\) 求 \(\displaystyle f^\prime(x)\)，并求 \(\displaystyle f(x)\) 的极值。


（16）（本题满分 10 分）设函数 \(\displaystyle f(u,\allowbreak v)\) 具有 2 阶连续偏导数，函数 \(\displaystyle g(x,\allowbreak y)=xy-f(x+y,\allowbreak x-y)\)。求 \(\displaystyle \dfrac{\displaystyle \partial^2g}{\displaystyle \partial x^2}+\dfrac{\displaystyle \partial^2g}{\displaystyle \partial x\partial y}+\dfrac{\displaystyle \partial^2g}{\displaystyle \partial y^2}\)。


（17）（本题满分 10 分）设函数 \(\displaystyle y(x)\) 是微分方程 \(\displaystyle y^\prime-xy=\dfrac{\displaystyle 1}{\displaystyle 2\sqrt x}e^{x^2/2}\) 满足条件 \(\displaystyle y(1)=\sqrt e\) 的特解。


（I）求 \(\displaystyle y(x)\)；


（II）设平面区域 \(\displaystyle D=\{(x,\allowbreak y)\mid1\leq x\leq2,\allowbreak 0\leq y\leq y(x)\}\)，求 \(\displaystyle D\) 绕 \(\displaystyle x\) 轴旋转所得旋转体的体积。


\clearpage

\subsection*{原 PDF 第 3 页}

（18）（本题满分 10 分）求曲线 \(\displaystyle y=e^{-x}\sin x\;(x\geq0)\) 与 \(\displaystyle x\) 轴之间图形的面积。


（19）（本题满分 10 分）设 \(\displaystyle a_n=\displaystyle \int_0^1x^n\sqrt{1-x^2}\,\mathrm dx\;(n=0,\allowbreak 1,\allowbreak 2,\allowbreak \cdots)\)。


（I）证明数列 \(\displaystyle \{a_n\}\) 单调递减，且 \(\displaystyle a_n=\dfrac{\displaystyle n-1}{\displaystyle n+2}a_{n-2}\;(n=2,\allowbreak 3,\allowbreak \cdots)\)；


（II）求 \(\displaystyle \lim_{n\to\infty}\dfrac{\displaystyle a_n}{\displaystyle a_{n-1}}\)。


（20）（本题满分 11 分）已知向量组 I：\(\displaystyle \boldsymbol\alpha_1=\begin{pmatrix}1\\1\\4\end{pmatrix},\allowbreak \ \boldsymbol\alpha_2=\begin{pmatrix}1\\0\\4\end{pmatrix},\allowbreak \ \boldsymbol\alpha_3=\begin{pmatrix}1\\2\\a^2+3\end{pmatrix}\) 与 II：\(\displaystyle \boldsymbol\beta_1=\begin{pmatrix}1\\1\\a+3\end{pmatrix},\allowbreak \ \boldsymbol\beta_2=\begin{pmatrix}0\\2\\1-a\end{pmatrix},\allowbreak \ \boldsymbol\beta_3=\begin{pmatrix}1\\3\\a^2+3\end{pmatrix}\)。若向量组 I 与 II 等价，求 \(\displaystyle a\) 的取值，并将 \(\displaystyle \boldsymbol\beta_3\) 用 \(\displaystyle \boldsymbol\alpha_1,\allowbreak \boldsymbol\alpha_2,\allowbreak \boldsymbol\alpha_3\) 线性表示。


\clearpage

\subsection*{原 PDF 第 4 页}

（21）（本题满分 11 分）已知矩阵 \(\displaystyle \boldsymbol A=\begin{pmatrix}-2&-2&1\\2&x&-2\\0&0&-2\end{pmatrix}\) 与 \(\displaystyle \boldsymbol B=\begin{pmatrix}2&1&0\\0&-1&0\\0&0&y\end{pmatrix}\) 相似。


（I）求 \(\displaystyle x,\allowbreak y\)；（II）求可逆矩阵 \(\displaystyle \boldsymbol P\)，使得 \(\displaystyle \boldsymbol P^{-1}\boldsymbol A\boldsymbol P=\boldsymbol B\)。


（22）（本题满分 11 分）设随机变量 \(\displaystyle X\) 与 \(\displaystyle Y\) 相互独立，\(\displaystyle X\) 服从参数为 1 的指数分布，\(\displaystyle Y\) 的概率分布为 \(\displaystyle P\{Y=-1\}=p,\allowbreak \ P\{Y=1\}=1-p\;(0<p<1)\)。令 \(\displaystyle Z=XY\)。


（I）求 \(\displaystyle Z\) 的概率密度；（II）\(\displaystyle p\) 为何值时，\(\displaystyle X\) 与 \(\displaystyle Z\) 不相关；（III）\(\displaystyle X\) 与 \(\displaystyle Z\) 是否相互独立？


（23）（本题满分 11 分）设总体 \(\displaystyle X\) 的概率密度为


\[\fitmath{\displaystyle f(x;\sigma^2)=\begin{cases}\dfrac{\displaystyle A}{\displaystyle \sigma} e^{-\dfrac{\displaystyle (x-\mu)^2}{\displaystyle 2\sigma^2}},&x\geq\mu,\\0,&x<\mu,\end{cases}}\]


其中 \(\displaystyle \mu\) 是已知参数，\(\displaystyle \sigma>0\) 是未知参数，\(\displaystyle A\) 是常数。\(\displaystyle X_1,\allowbreak X_2,\allowbreak \cdots,\allowbreak X_n\) 是来自总体 \(\displaystyle X\) 的简单随机样本。


（I）求 \(\displaystyle A\)；（II）求 \(\displaystyle \sigma^2\) 的最大似然估计量。


\clearpage
\end{document}
